\section{Mean-support growth along an orbit}\label{sec:mean}
For a planar convex body $P$, let $h_P(u)=\max_{z\in P}u\cdot z$ and
\[
 \mw(P)=\frac1{2\pi}\int_0^{2\pi}h_P(\cos\theta,\sin\theta)\,d\theta.
\]
For polygons this is perimeter divided by $2\pi$, by Cauchy's perimeter
formula. Define
\[
 \mathcal G_\alpha(P)=\conv(P\cup R_\alpha P),\qquad
 D_\alpha(P)=\mw(\mathcal G_\alpha(P))-\mw(P).
\]
The functional is monotone under inclusion and invariant under rotation.
For a chain \eqref{eq:chain},
\begin{equation}\label{eq:mean-budget}
 \sum_{t<N}D_\alpha(P_t)
 \le\mw(P_N)-\mw(P_0)\le\sqrt2.
\end{equation}
The last bound uses $P_N\subset[-1,1]^2\subset\sqrt2\B^2$.
It is deliberately larger than the exact endpoint allowance.

\begin{lemma}[Accumulation without squaring]\label{lem:accumulation}
For any integer $q\ge2$ and $P$,
\begin{equation}\label{eq:orbit-sum}
 \mw\!\left(\conv\bigcup_{k=0}^{q-1}R_\alpha^k P\right)-\mw(P)
                          \le(q-1)D_\alpha(P).
\end{equation}
\end{lemma}
\begin{proof}
Write $h_k(\theta)=h_P(\theta-2\pi k\alpha)$. Pointwise,
$\max_{k<q}h_k-h_0\le\sum_{k=1}^{q-1}(h_k-h_{k-1})_+$.
Each summand has the same integral as $(h_1-h_0)_+$, whose normalized
integral is $D_\alpha(P)$. Integrate the inequality.
\end{proof}

\begin{lemma}[Inscribed-polygon deficit]\label{lem:perimeter}
Suppose $a\B^2\subset P$, and $P$ has at most $M\ge3$ vertices.
Put $r=\max_{z\in P}\norm z_2$. Then
\begin{equation}\label{eq:perimeter}
                 \mw(P)\le r\frac M\pi\sin\frac\pi M.
\end{equation}
\end{lemma}
\begin{proof}
Push each vertex radially to the circle of radius $r$. The resulting
convex hull still contains zero and hence contains the old vertices and
$P$. If its successive angular gaps are $\theta_1,\ldots,\theta_m$,
$m\le M$, its perimeter is $2r\sum_i\sin(\theta_i/2)$.
Concavity of sine on $[0,\pi]$, $\sum_i\theta_i=2\pi$, and monotonicity
of $m\sin(\pi/m)$ give at most $2rM\sin(\pi/M)$.
Monotonicity of mean support and Cauchy's formula prove the assertion.
The radial hull contains zero because any positive affine dependence
of the old vertices giving zero can be reweighted by their radii.
\end{proof}

\begin{lemma}[A denominator gives an orbit net]\label{lem:net}
If $(p,q)=1$ and $|\alpha-p/q|<q^{-2}$, the first $q$ points of the
rotation orbit of any unit vector form an angular $3\pi/q$-net.
\end{lemma}
\begin{proof}
The multiples of $p/q$ give the uniform grid of mesh $1/q$ turns.
The corresponding first $q$ multiples of $\alpha$ differ by less than
$1/q$ turns. Every point is within $1/(2q)$ turns of the grid, proving
the bound after multiplying by $2\pi$.
\end{proof}

\begin{theorem}[One-step quantitative mean-support gain]
\label{thm:mean-growth}
Suppose $a\B^2\subset P$, $|\ext P|\le M$, $M\ge3$, and a reduced
fraction $p/q$ satisfies $q\ge8M$ and $|\alpha-p/q|<q^{-2}$. Then
\begin{equation}\label{eq:gain-q}
                        D_\alpha(P)\ge\frac{3a}{4M^2q}.
\end{equation}
\end{theorem}
\begin{proof}
Take a vertex of maximum radius $r$. Its first $q$ rotations and
Lemma~\ref{lem:net} imply
\[
 \mw\!\left(\conv\bigcup_{k<q}R_\alpha^k P\right)
                              \ge r\cos(3\pi/q).
\]
Subtract \eqref{eq:perimeter}. Since $M\ge3$ and $q\ge8M$,
Taylor's inequalities for sine and cosine give
\begin{align*}
 \cos(3\pi/q)-\frac M\pi\sin(\pi/M)
 &\ge \frac{\pi^2}{M^2}
       \left(\frac16-\frac{\pi^2}{120M^2}-\frac{9M^2}{2q^2}\right)\\
 &\ge\frac9{M^2}\left(\frac16-\frac{10}{1080}-\frac9{128}\right)
   =\frac{301}{384M^2}>\frac3{4M^2}.
\end{align*}
The inequalities use only $9<\pi^2<10$. Apply
Lemma~\ref{lem:accumulation}, use $r\ge a$ and $q-1<q$.
The gain is linear in the integrated support deficit; an area cap is
not introduced.
\end{proof}

For an irrational $\alpha\in(0,1)$, let $p_k/q_k$ be its continued-fraction
convergents, ignoring the initial repeated denominator one, and define
\begin{equation}\label{eq:denominator-function}
 Q_\alpha(M)=\min\{q_k:q_k\ge8M\}.
\end{equation}
We use the elementary relations
\begin{equation}\label{eq:cf}
 (p_k,q_k)=1,\quad
 |\alpha-p_k/q_k|<\frac1{q_kq_{k+1}}\le\frac1{q_k^2},\quad
 q_{k+1}=a_{k+1}q_k+q_{k-1}.
\end{equation}
They follow by writing $\alpha$ with its remaining continued-fraction
complete quotient and using
$p_kq_{k-1}-p_{k-1}q_k=(-1)^{k-1}$; see also \cite{Khinchin}.
In particular, if $a_k\le A$ for all $k\ge1$,
$Q_\alpha(M)<8(A+1)M$.

\begin{theorem}[Profile budget and bounded-type width]
\label{thm:profiles}
Every exact hidden-state realization of \eqref{eq:array} satisfies
\begin{equation}\label{eq:general-profile}
 \sum_{t<N}\frac1{2^{2K_t}Q_\alpha(2^{K_t})}
                                           \le\frac{40\sqrt2}{3}.
\end{equation}
If $a_k\le A$, then
\begin{equation}\label{eq:bounded-profile}
              \sum_{t<N}2^{-3K_t}\le160(A+1).
\end{equation}
For vector-state realizations the same statements hold with $2^{K_t}$
replaced by $K_t$. Consequently
$W_N\ge\frac13\log_2(N/[160(A+1)])$ and
$V_N\ge(N/[160(A+1)])^{1/3}$.
\end{theorem}
\begin{proof}
Use Lemma~\ref{lem:reachable} and Theorem~\ref{thm:mean-growth} with
$a=1/10$, $M=2^{K_t}$, $q=Q_\alpha(M)$ in
\eqref{eq:mean-budget}. This gives \eqref{eq:general-profile}.
Under the bounded-type assumption each gain is at least
$3/[320(A+1)M^3]$. The endpoint allowance gives
\[
 \sum_{t<N}M^{-3}\le\frac{320(A+1)\sqrt2}{3}<160(A+1).
\]
In the vector-state class, use the other vertex bound in
Lemma~\ref{lem:reachable}. A peak bound supplies the stated consequences.
\end{proof}

For the golden angle all partial quotients equal one. The explicit hidden
lower bound exceeds three when $N>320\,2^9=163840$.
The earlier area estimate exceeded three only after
$24000000\,2^{18}=6291456000000$ commands \cite{GTF32}.
The new constant is still not a practically sharp finite-horizon
classification. Its significance is the improved exponent in the profile
budget and its matching vector-state construction below.

The same mean-support proof works in any dimension for a specified
orthogonal command family once one supplies an orbit-covering estimate
and an inscribed-polytope mean-support deficit. No higher-dimensional
numerical exponent is asserted here. Unlike a volume argument, the
potential is not tied to a determinant formula. A strictly contractive
rotation instead has a finite invariant polygon; the divergent conclusion
cannot simply be retained under contraction.
